\documentclass[11pt]{article}

\usepackage{jcd-macros}
\usepackage[numbers,square]{natbib}
\bluehyperref  % Dark blue hyperref
\usemedgeometry  % Use medium geometry (3cm left/right margin, 2.54 top/bottom)

\begin{document}

\title{26-11 Questions on the Inner Product Laplacian Revisited}

\maketitle

\section{Some minor notes on discussion}

\paragraph{Responsible Parties:} John, Tselil

\begin{enumerate}[1.]
\item Our goals on this one: for the first part of the question, try to
  demonstrate that with random signs $S \sim \uniform(\diag(\{\pm
  1\}^n))$ in the edge-adjacency matrix $B$, $M_V^{-1/2} S B M_E B^T S^T
  M_V^{-1/2}$ concentrates quickly (in operator norm, say) to
  $M_V^{-1/2} B M_E B^T M_V^{-1/2}$, so there's little randomness.
\item Give examples of the extraordinary fragility of eigenvectors without
  separation.
  %
  A trivial example of this comes from the matrices
  \begin{align*}
    \left[\begin{matrix} 1 - \epsilon & 0 \\ 0 & 1 + \epsilon
      \end{matrix} \right]
    ~~ \mbox{and} ~~
    \left[\begin{matrix} 1 & \epsilon \\ \epsilon & 1 + 2 \epsilon
        \end{matrix}\right].
  \end{align*}
  The first has $e_1$ and $e_2$ as eigenvectors, while the second has
  very different (unique) eigenvectors.
\end{enumerate}

\end{document}
